% problem-000387 5 反应动力学与催化
\section*{5. 反应动力学与催化}
$
\mathrm{CH} _ {3} \mathrm{COCH} _ {3} \xrightarrow {k _ {1}} \mathrm{CH} _ {3} \cdot + \mathrm{CH} _ {3} \mathrm{CO} \cdot
$

$
E _ {a} (1) = 3 5 1 \mathrm{kJ} \mathrm{mol} ^ {- 1}
$

$
\mathrm{CH} _ {3} \mathrm{CO} \cdot \xrightarrow {k _ {2}} \mathrm{CH} _ {3} \cdot + \mathrm{CO}
$

$
E _ {a} (2) = 4 2 \mathrm{kJ} \mathrm{mol} ^ {- 1}
$

$
\mathrm{CH} _ {3} \cdot + \mathrm{CH} _ {3} \mathrm{COCH} _ {3} \xrightarrow {k _ {3}} \mathrm{CH} _ {4} + \mathrm{CH} _ {3} \mathrm{COCH} _ {2} \cdot
$

$
E _ {a} (3) = 6 3 \mathrm{kJ} \mathrm{mol} ^ {- 1}
$

$
\mathrm{CH} _ {3} \mathrm{COCH} _ {2} \cdot \xrightarrow {k _ {4}} \mathrm{CH} _ {3} \cdot + \mathrm{CH} _ {2} \mathrm{CO}\tag{4}
$

$
E _ {a} (4) = 2 0 0 \mathrm{kJ} \mathrm{mol} ^ {- 1}
$

$
\mathrm{CH} _ {3} \cdot + \mathrm{CH} _ {3} \mathrm{COCH} _ {2} \cdot \xrightarrow {k _ {5}} \mathrm{C} _ {2} \mathrm{H} _ {5} \mathrm{COCH} _ {3}
$

$
E _ {a} (5) = 2 1 \mathrm{kJ} \mathrm{mol} ^ {- 1}
$

\subsection*{5-1}
上述历程是根据两个平⾏反应⽽提出的，请写出反应⽅程式。

\subsection*{5-2}
实验测得1000K下热分解反应对丙酮为⼀级反应，试推出反应的速率⽅程。提⽰：①消耗丙酮的主要基元步骤是反应(3)；②可忽略各基元反应指前因⼦的影响。

\subsection*{5-3}
根据5-2推出的速率⽅程，计算表观活化能。

\subsection*{5-4}
分解产物中 $\mathrm { C H _ { 4 } }$ 可以和⽔蒸⽓反应制备H_{2}：

$
\mathrm{CH} _ {4} (\mathrm{g}) + \mathrm{H} _ {2} \mathrm{O} (\mathrm{g}) \rightarrow 3 \mathrm{H} _ {2} (\mathrm{g}) + \mathrm{CO} (\mathrm{g})
$

试根据下表数据计算该反应在298K时 $\Delta _ { r } G _ { m } ^ { \ominus }$ 和平衡常数 $K ^ { \ominus }$ ，并指出平衡常数随温度的变化趋势。

相关物质热⼒学数据(298K)

\textit{（原卷此处含表格/仪器试剂表，详见 PDF 或 MD 源文件。）}
